\section{The Compounding Advantage}
\label{sec:learning}

Model progress expands what Fusion can draw upon. Allocation progress determines
how much of that supply becomes useful system capability. The latter is where
Fusion develops its own compounding advantage. An asymmetric model pool can be
copied and a routing graph can be reconstructed; an empirical understanding of
how capabilities interact with real tasks---and a system that improves that
understanding through use---is harder to reproduce.

\subsection{A living map of model capability}

Leaderboards describe average model performance on predetermined tasks. Fusion
needs a different object: a conditional capability map. It asks how likely a
model is to contribute useful evidence for a particular kind of task, under a
particular role, at a particular workflow state, given the evidence already
available.

The map includes capability and operating characteristics together: model and
provider version, task dimensions, role, context, output validity, latency,
cost, observed redundancy, error correlation, synthesizer use, and downstream
outcome. It is versioned because models change. It is uncertainty-aware because
new models and rare tasks begin with limited evidence. And it measures
interaction, because the best standalone answerer is not necessarily the most
valuable next contributor.

This systematic evaluation is what makes heterogeneous allocation reliable.
Without it, asymmetry is merely an intuition that cheaper models might help.
With it, the system can predict where a capability is complementary, where it
is redundant, and where the risk requires escalation.

\subsection{The allocation data flywheel}

Every Fusion task can produce a structured learning record: the state the
system observed, the uncertainty it identified, the call it admitted, the role
and model it selected, the evidence returned, the decision committed, the
resources consumed, and the eventual outcome.

These records improve three layers at once:

\begin{itemize}
  \item \textbf{task understanding}: a richer representation of what real users
  are trying to accomplish and where their workflows become uncertain;
  \item \textbf{capability understanding}: empirical knowledge of what each
  model contributes under real conditions; and
  \item \textbf{allocation policy}: better admission, role assignment,
  escalation, synthesis, and stopping decisions.
\end{itemize}

More use creates a wider view of task demand. Better task understanding creates
better allocation. Better allocation creates more informative observations
about model capability. The result is a flywheel that is grounded not only in
benchmark labels, but in the distribution of real tasks and feedback that the
system encounters.

This flywheel is difficult to reproduce from the surface behavior of Fusion.
It depends on the joint history of tasks, model versions, allocations,
counterfactual probes, and outcomes. The visible architecture is modular; the
accumulated model of conditional contribution is proprietary and compounding.

\subsection{From designed allocation to learned allocation}

Fusion currently combines structured task perception with explicit policy
boundaries. This makes allocation interpretable while the learning base grows.
The same architecture exposes a stable end-to-end learning interface:

\begin{align}
 &(\text{task state},\text{available evidence},\text{model supply}) \notag\\
 &\hspace{2.0em}\longrightarrow
 (\text{next source},\text{role},\text{budget},\text{stop}).
\end{align}

Fumo Lab has begun training models that map user input and evolving task state
directly to allocation decisions. Producer behavior supplies model feedback:
which sources generated novel, valid, or decision-changing evidence.
Synthesizer behavior and downstream outcomes supply system feedback: which
evidence was used, which commitment succeeded, and where the workflow failed.
Together, they create a path toward reinforcement learning from model and
environment feedback for the allocation policy itself.

Learning from production traces requires discipline. The system only observes
the outcomes of calls it chose to make, so naive training would reinforce its
existing blind spots. Controlled exploration, shadow calls, candidate
ablations, offline replay, and local credit assignment are necessary to
estimate what would have happened under another allocation. Privacy,
provenance, evaluation, and rollback must be properties of the learning loop,
not afterthoughts.

\subsection{Bounded system-level recursive improvement}

Fusion improves through two coupled streams. New and better models expand the
supply of intelligence. Better task understanding and allocation extract more
value from that supply. Improvements in either stream change the data available
to the other: stronger producers reveal new possibilities, while a stronger
allocator discovers where each producer matters.

This is recursive improvement in a bounded, system-level sense. Fusion uses
observations of its own decisions and outcomes to improve how future
intelligence is allocated. It does not imply unrestricted self-modification or
an unobservable loop. Policies remain versioned, evaluated, constrained, and
reversible.

\principle{Fusion improves when models improve---and improves again when it
learns how to use them.}
